% =====================================================================
%  Kapitel 1 – Einleitung (~1 Seite)
% =====================================================================
\chapter{Einleitung}
\label{kap:einleitung}

Stress betrifft sehr viele Menschen, sei es durch Arbeit, Studium oder belastende
Lebensumstände. Eine einfache Möglichkeit, Anspannung zu begegnen, sind angeleitete
Atemübungen mit einem ruhigen, gleichmäßigen Rhythmus aus Ein- und Ausatmen. Soziale Roboter
können solche Übungen nicht nur über Bild und Ton anleiten, sondern den Atemrhythmus auch
körperlich spürbar machen. Der an der Yale University entwickelte Roboter Ommie gibt Einatmen,
Halten und Ausatmen über haptische Interaktion und Audio vor~\cite{matheus2022ommie}. Im Labor
für Kommunikation der \ac{BHT} wurde dieses Prinzip mit Neffy aufgegriffen: In einer
Pilotstudie mit Neffy~1.0 sanken die Angstwerte der Teilnehmenden nach der Übung
signifikant~\cite{thomas2025neffy}. In einer Studie mit 14 Geflüchteten aus der Ukraine
empfanden die Teilnehmenden ihren Stress nach einer Übung mit Neffy~2.0 als deutlich geringer
als nach einer reinen Audioanleitung~\cite{buchem2026neffy2}.

Die Zielgruppe von Neffy ist bewusst weit gefasst. Denkbar ist der Einsatz im Alltag ebenso
wie bei Kindern oder bei besonders belasteten Gruppen wie in der Studie zu Neffy~2.0. Ein
Anliegen des Gesamtprojekts ist es, herauszufinden, in welchen Anwendungsfällen ein solcher
Roboter am meisten Sinn ergibt; möglich sind später auch verschiedene Varianten für
unterschiedliche Einsatzbereiche. Neffy~3.0 ist in diesem Sinne ein Zwischenschritt in einem
fortlaufenden Entwicklungsprozess.

Neffy~1.0 und 2.0 geben den Atemrhythmus vor, nehmen ihr Gegenüber aber nicht wahr. Mit
Neffy~3.0 soll sich das ändern: Der Roboter soll erkennen, dass jemand vor ihm sitzt, die
Atmung der Nutzer:innen erfassen und darauf mit einer passenden Übung reagieren. Ursprünglich
war dafür eine vollautomatische Erfassung vorgesehen, bei der die Atmung allein über die
aufliegenden Hände erkannt wird. Wie Kapitel~\ref{kap:diskussion} zeigt, ließ sich das mit dem
erprobten Aufbau nicht zuverlässig umsetzen. Die umgesetzte Lösung erfasst stattdessen den
Atemtakt, den die Nutzer:innen angeben, indem sie beim Einatmen aktiv in die Schultertaschen
drücken; sie setzt also ihre bewusste Mitwirkung voraus.

Dieser Bericht beschreibt den Beitrag des Teams Sensorik zu Neffy~3.0. Das Team war zunächst
für die Anwesenheits- und die Atemerkennung zuständig und übernahm im Verlauf des Projekts
zusätzlich die Konstruktion und Fertigung des neuen Gehäuses. Kapitel~\ref{kap:ziel} legt die
Ziele, die Anforderungen und die Organisation des Projekts dar. Kapitel~\ref{kap:stand}
ordnet die gewählten Verfahren in den Stand der Technik ein und vergleicht Neffy~3.0 mit
seinen Vorgängern. Kapitel~\ref{kap:methoden} beschreibt Systemarchitektur, Sensorik,
Gehäuse und Firmware, Kapitel~\ref{kap:diskussion} die Ergebnisse, die Messwerte und die
Bewertung gegen die Anforderungen. Kapitel~\ref{kap:fazit} fasst die Arbeit zusammen und
gibt einen Ausblick, Kapitel~\ref{kap:einschaetzung} enthält unsere persönliche Einschätzung
des Praktikums.
